\documentclass[12pt,a4paper]{exam}
\usepackage[utf8]{inputenc}
\usepackage{amsmath,amssymb,amsthm}
\usepackage[margin=1in]{geometry}
\usepackage{graphicx}
\usepackage[colorlinks=true]{hyperref}
\pagestyle{headandfoot}
\firstpageheader{MIT OpenCourseWare}{}{\textbf{EXTREMELY HARD -- MIT LEVEL}}
\runningheader{MIT OpenCourseWare}{}{\textbf{EXTREMELY HARD -- MIT LEVEL}}
\firstpagefooter{}{Page \thepage}{}
\runningfooter{}{Page \thepage}{}

\begin{document}

\begin{center}
{\LARGE \textbf{MIT OpenCourseWare}}\\8.02SC\\Electricity and Magnetism
\vspace{0.5cm}
{8.02SC} --- {Electricity and Magnetism}
\vspace{0.3cm}
May 2026
\vspace{0.3cm}
\textbf{EXTREMELY HARD -- MIT LEVEL}
\end{center}

\begin{center}
\textbf{Time Limit: 3 hours}
\end{center}

\begin{center}
\textbf{Total Marks: 100}
\end{center}

\vspace{0.5cm}

\begin{questions}

\question[2]
Construct an original proof-level problem involving Electrostatics for 8.02SC (Electricity and Magnetism) and provide a complete solution. Your problem should be at the level of a graduate qualifying exam.

\vspace{0.3cm}

\question[3]
Analyze the limitations of standard approaches to Magnetism when applied to edge cases in 8.02SC. Construct explicit counterexamples showing where intuitive reasoning fails.

\vspace{0.3cm}

\question[3]
Construct an original proof-level problem involving Maxwell's Equations for 8.02SC (Electricity and Magnetism) and provide a complete solution. Your problem should be at the level of a graduate qualifying exam.

\vspace{0.3cm}

\question[3]
Construct an original proof-level problem involving EM Waves for 8.02SC (Electricity and Magnetism) and provide a complete solution. Your problem should be at the level of a graduate qualifying exam.

\vspace{0.3cm}

\question[4]
Construct an original proof-level problem involving Circuit Theory for 8.02SC (Electricity and Magnetism) and provide a complete solution. Your problem should be at the level of a graduate qualifying exam.

\vspace{0.3cm}

\question[2]
Construct an original proof-level problem involving Potential Theory for 8.02SC (Electricity and Magnetism) and provide a complete solution. Your problem should be at the level of a graduate qualifying exam.

\vspace{0.3cm}

\question[5]
Prove the fundamental theorem concerning Electrostatics in the context of 8.02SC (Electricity and Magnetism). State the theorem precisely, provide a complete proof, and discuss three important corollaries.

\vspace{0.3cm}

\question[2]
Analyze the limitations of standard approaches to Magnetism when applied to edge cases in 8.02SC. Construct explicit counterexamples showing where intuitive reasoning fails.

\vspace{0.3cm}

\question[3]
Prove the fundamental theorem concerning Maxwell's Equations in the context of 8.02SC (Electricity and Magnetism). State the theorem precisely, provide a complete proof, and discuss three important corollaries.

\vspace{0.3cm}

\question[5]
Prove the fundamental theorem concerning EM Waves in the context of 8.02SC (Electricity and Magnetism). State the theorem precisely, provide a complete proof, and discuss three important corollaries.

\vspace{0.3cm}

\question[4]
Construct an original proof-level problem involving Circuit Theory for 8.02SC (Electricity and Magnetism) and provide a complete solution. Your problem should be at the level of a graduate qualifying exam.

\vspace{0.3cm}

\question[6]
Construct an original proof-level problem involving Potential Theory for 8.02SC (Electricity and Magnetism) and provide a complete solution. Your problem should be at the level of a graduate qualifying exam.

\vspace{0.3cm}

\question[5]
Prove the fundamental theorem concerning Electrostatics in the context of 8.02SC (Electricity and Magnetism). State the theorem precisely, provide a complete proof, and discuss three important corollaries.

\vspace{0.3cm}

\question[4]
Prove the fundamental theorem concerning Magnetism in the context of 8.02SC (Electricity and Magnetism). State the theorem precisely, provide a complete proof, and discuss three important corollaries.

\vspace{0.3cm}

\question[4]
Construct an original proof-level problem involving Maxwell's Equations for 8.02SC (Electricity and Magnetism) and provide a complete solution. Your problem should be at the level of a graduate qualifying exam.

\vspace{0.3cm}

\question[6]
Prove the fundamental theorem concerning EM Waves in the context of 8.02SC (Electricity and Magnetism). State the theorem precisely, provide a complete proof, and discuss three important corollaries.

\vspace{0.3cm}

\question[4]
Analyze the limitations of standard approaches to Circuit Theory when applied to edge cases in 8.02SC. Construct explicit counterexamples showing where intuitive reasoning fails.

\vspace{0.3cm}

\question[4]
Construct an original proof-level problem involving Potential Theory for 8.02SC (Electricity and Magnetism) and provide a complete solution. Your problem should be at the level of a graduate qualifying exam.

\vspace{0.3cm}

\question[5]
Prove the fundamental theorem concerning Electrostatics in the context of 8.02SC (Electricity and Magnetism). State the theorem precisely, provide a complete proof, and discuss three important corollaries.

\vspace{0.3cm}

\question[4]
Prove the fundamental theorem concerning Magnetism in the context of 8.02SC (Electricity and Magnetism). State the theorem precisely, provide a complete proof, and discuss three important corollaries.

\vspace{0.3cm}

\question[5]
Analyze the limitations of standard approaches to Maxwell's Equations when applied to edge cases in 8.02SC. Construct explicit counterexamples showing where intuitive reasoning fails.

\vspace{0.3cm}

\question[4]
Analyze the limitations of standard approaches to EM Waves when applied to edge cases in 8.02SC. Construct explicit counterexamples showing where intuitive reasoning fails.

\vspace{0.3cm}

\question[4]
Analyze the limitations of standard approaches to Circuit Theory when applied to edge cases in 8.02SC. Construct explicit counterexamples showing where intuitive reasoning fails.

\vspace{0.3cm}

\question[6]
Prove the fundamental theorem concerning Potential Theory in the context of 8.02SC (Electricity and Magnetism). State the theorem precisely, provide a complete proof, and discuss three important corollaries.

\vspace{0.3cm}

\question[3]
Prove the fundamental theorem concerning Electrostatics in the context of 8.02SC (Electricity and Magnetism). State the theorem precisely, provide a complete proof, and discuss three important corollaries.

\vspace{0.3cm}

\end{questions}

\newpage
\section*{Solutions}

\textbf{Solution 1:}

Solution approach: First recall the definitions and key theorems related to Electrostatics from 8.02SC. Then apply the appropriate techniques to construct a rigorous argument. The solution must be self-contained and reference only the material from the course.

\vspace{0.5cm}

\textbf{Solution 2:}

The edge case analysis reveals the subtle assumptions underlying standard treatments of Magnetism. By constructing explicit counterexamples, we see that the usual hypotheses cannot be weakened without loss of the theorem's conclusion.

\vspace{0.5cm}

\textbf{Solution 3:}

Solution approach: First recall the definitions and key theorems related to Maxwell's Equations from 8.02SC. Then apply the appropriate techniques to construct a rigorous argument. The solution must be self-contained and reference only the material from the course.

\vspace{0.5cm}

\textbf{Solution 4:}

Solution approach: First recall the definitions and key theorems related to EM Waves from 8.02SC. Then apply the appropriate techniques to construct a rigorous argument. The solution must be self-contained and reference only the material from the course.

\vspace{0.5cm}

\textbf{Solution 5:}

Solution approach: First recall the definitions and key theorems related to Circuit Theory from 8.02SC. Then apply the appropriate techniques to construct a rigorous argument. The solution must be self-contained and reference only the material from the course.

\vspace{0.5cm}

\textbf{Solution 6:}

Solution approach: First recall the definitions and key theorems related to Potential Theory from 8.02SC. Then apply the appropriate techniques to construct a rigorous argument. The solution must be self-contained and reference only the material from the course.

\vspace{0.5cm}

\textbf{Solution 7:}

The proof follows from the definitions and properties established in 8.02SC. The key insight is to apply the relevant theorem and verify the hypotheses hold. Detailed solution depends on the specific formulation of the theorem.

\vspace{0.5cm}

\textbf{Solution 8:}

The edge case analysis reveals the subtle assumptions underlying standard treatments of Magnetism. By constructing explicit counterexamples, we see that the usual hypotheses cannot be weakened without loss of the theorem's conclusion.

\vspace{0.5cm}

\textbf{Solution 9:}

The proof follows from the definitions and properties established in 8.02SC. The key insight is to apply the relevant theorem and verify the hypotheses hold. Detailed solution depends on the specific formulation of the theorem.

\vspace{0.5cm}

\textbf{Solution 10:}

The proof follows from the definitions and properties established in 8.02SC. The key insight is to apply the relevant theorem and verify the hypotheses hold. Detailed solution depends on the specific formulation of the theorem.

\vspace{0.5cm}

\textbf{Solution 11:}

Solution approach: First recall the definitions and key theorems related to Circuit Theory from 8.02SC. Then apply the appropriate techniques to construct a rigorous argument. The solution must be self-contained and reference only the material from the course.

\vspace{0.5cm}

\textbf{Solution 12:}

Solution approach: First recall the definitions and key theorems related to Potential Theory from 8.02SC. Then apply the appropriate techniques to construct a rigorous argument. The solution must be self-contained and reference only the material from the course.

\vspace{0.5cm}

\textbf{Solution 13:}

The proof follows from the definitions and properties established in 8.02SC. The key insight is to apply the relevant theorem and verify the hypotheses hold. Detailed solution depends on the specific formulation of the theorem.

\vspace{0.5cm}

\textbf{Solution 14:}

The proof follows from the definitions and properties established in 8.02SC. The key insight is to apply the relevant theorem and verify the hypotheses hold. Detailed solution depends on the specific formulation of the theorem.

\vspace{0.5cm}

\textbf{Solution 15:}

Solution approach: First recall the definitions and key theorems related to Maxwell's Equations from 8.02SC. Then apply the appropriate techniques to construct a rigorous argument. The solution must be self-contained and reference only the material from the course.

\vspace{0.5cm}

\textbf{Solution 16:}

The proof follows from the definitions and properties established in 8.02SC. The key insight is to apply the relevant theorem and verify the hypotheses hold. Detailed solution depends on the specific formulation of the theorem.

\vspace{0.5cm}

\textbf{Solution 17:}

The edge case analysis reveals the subtle assumptions underlying standard treatments of Circuit Theory. By constructing explicit counterexamples, we see that the usual hypotheses cannot be weakened without loss of the theorem's conclusion.

\vspace{0.5cm}

\textbf{Solution 18:}

Solution approach: First recall the definitions and key theorems related to Potential Theory from 8.02SC. Then apply the appropriate techniques to construct a rigorous argument. The solution must be self-contained and reference only the material from the course.

\vspace{0.5cm}

\textbf{Solution 19:}

The proof follows from the definitions and properties established in 8.02SC. The key insight is to apply the relevant theorem and verify the hypotheses hold. Detailed solution depends on the specific formulation of the theorem.

\vspace{0.5cm}

\textbf{Solution 20:}

The proof follows from the definitions and properties established in 8.02SC. The key insight is to apply the relevant theorem and verify the hypotheses hold. Detailed solution depends on the specific formulation of the theorem.

\vspace{0.5cm}

\textbf{Solution 21:}

The edge case analysis reveals the subtle assumptions underlying standard treatments of Maxwell's Equations. By constructing explicit counterexamples, we see that the usual hypotheses cannot be weakened without loss of the theorem's conclusion.

\vspace{0.5cm}

\textbf{Solution 22:}

The edge case analysis reveals the subtle assumptions underlying standard treatments of EM Waves. By constructing explicit counterexamples, we see that the usual hypotheses cannot be weakened without loss of the theorem's conclusion.

\vspace{0.5cm}

\textbf{Solution 23:}

The edge case analysis reveals the subtle assumptions underlying standard treatments of Circuit Theory. By constructing explicit counterexamples, we see that the usual hypotheses cannot be weakened without loss of the theorem's conclusion.

\vspace{0.5cm}

\textbf{Solution 24:}

The proof follows from the definitions and properties established in 8.02SC. The key insight is to apply the relevant theorem and verify the hypotheses hold. Detailed solution depends on the specific formulation of the theorem.

\vspace{0.5cm}

\textbf{Solution 25:}

The proof follows from the definitions and properties established in 8.02SC. The key insight is to apply the relevant theorem and verify the hypotheses hold. Detailed solution depends on the specific formulation of the theorem.

\vspace{0.5cm}

\end{document}